\chapter{Mathematics Course}

\section{What The Math Is For}

The auto-aim system answers one question many times per second:

\begin{quote}
Given a camera image, the current gimbal yaw and pitch, and a moving target,
what absolute yaw and pitch should the gimbal move to, and is it safe to fire?
\end{quote}

The repository solves that question as a chain of smaller math problems:

\begin{enumerate}
  \item A detector reports a 2D rectangle around an armor plate.
  \item PnP estimates where that plate is in 3D camera coordinates.
  \item A frame transform moves the 3D point into the robot's odom frame.
  \item An Extended Kalman Filter estimates the enemy robot center, velocity,
  yaw, yaw rate, and armor radius.
  \item The aim planner predicts which of four armor faces should be hit at
  bullet impact time.
  \item The ballistic solver computes the pitch needed to compensate gravity.
  \item The fire gate checks whether the command is accurate enough to shoot.
\end{enumerate}

Each step has its own units. Mixing units is one of the fastest ways to break
this project. In this package, distances are meters unless a comment says
millimeters, angles are radians unless a parameter name ends with
\code{_deg}, and image coordinates are pixels.

\section{Math Notation For Beginners}

A scalar is one number, such as $3.0$ meters. A vector is a list of numbers,
such as a 3D position:

\[
\vect{p} =
\begin{bmatrix}
x \\ y \\ z
\end{bmatrix}.
\]

A matrix is a rectangular table of numbers. Multiplying a matrix by a vector
can rotate, scale, or mix the vector components. A rotation matrix $R$ and a
translation vector $\vect{t}$ transform a point from one frame to another:

\[
\vect{p}_{world} = R \vect{p}_{camera} + \vect{t}.
\]

The code often uses the same idea with quaternions instead of explicit
rotation matrices. A quaternion is another way to store a 3D rotation. You do
not need to manually multiply quaternions at first; you need to know that
\code{tf2::quatRotate(q, p)} rotates point $\vect{p}$ by rotation $q$.

\section{Coordinate Frames}

A coordinate frame defines what the axes mean. The same physical armor can
have different coordinates in different frames.

\begin{longtable}{p{0.24\textwidth}p{0.65\textwidth}}
\toprule
Frame & Meaning in this repository \\
\midrule
Camera optical & ROS camera convention: $+Z$ forward from the camera,
$+X$ right in the image, $+Y$ down in the image. PnP outputs points here. \\
Gimbal body & Mechanical convention: $+X$ forward, $+Y$ left, $+Z$ up. The
barrel offset is measured in this frame. \\
Odom / micro reference & Fixed startup reference used by the microcontroller.
The node publishes absolute yaw and pitch commands in this frame after sign
conversion. \\
\bottomrule
\end{longtable}

In \file{src/src/frame_transformer.cpp}, the current IMU yaw and pitch build a
rotation:

\[
q_{camera\rightarrow odom} = q_{gimbal} q_{body\rightarrow camera}.
\]

Then a PnP point is transformed as:

\[
\vect{p}_{odom} =
q_{camera\rightarrow odom} \vect{p}_{camera} +
\begin{bmatrix}0\\0\\h_{gimbal}\end{bmatrix}.
\]

The code form is:

\begin{lstlisting}[language=C++]
tf2::Vector3 p_odom = tf2::quatRotate(rotation_, p_cam) + translation_;
\end{lstlisting}

\subsection{Yaw And Pitch}

Yaw turns left/right around the vertical axis. Pitch turns up/down. The command
published on \topic{/tracker/cmd_gimbal} is not a small change. It is an
absolute destination:

\[
\text{command} =
\begin{cases}
\text{absolute yaw target in radians}\\
\text{absolute pitch target in radians}
\end{cases}
\]

The parameters \param{gimbal.yaw_sign} and \param{gimbal.pitch_sign} convert
between the code's internal convention and the microcontroller's convention.
They must be exactly $+1$ or $-1$.

\section{Camera Projection}

A camera maps a 3D point to a 2D pixel. Ignoring distortion for the moment:

\[
u = f_x \frac{X}{Z} + c_x,
\qquad
v = f_y \frac{Y}{Z} + c_y.
\]

Here $(X,Y,Z)$ is a point in camera optical coordinates, $(u,v)$ is the pixel,
$f_x$ and $f_y$ are focal lengths in pixels, and $(c_x,c_y)$ is the optical
center in pixels. These values come from \topic{/camera_info}.

The code also uses the inverse idea for the HUD overlay. If the aim ray has a
small yaw and pitch in the camera frame, the pixel is:

\[
u = c_x + f_x \tan(\theta_{yaw}),
\qquad
v = c_y - f_y \tan(\theta_{pitch}).
\]

This appears in \file{src/src/auto_aim_node.cpp} when publishing
\topic{/tracker/aim_pixels}.

\section{PnP: From 2D Rectangle To 3D Pose}

PnP means ``Perspective-n-Point''. It estimates the camera-frame pose of a
known 3D object from corresponding 3D object points and 2D image points.

This repository does not get exact armor corner points from the detector. It
gets a bounding box: center $(c_x,c_y)$, width $w$, and height $h$. The PnP
solver builds a synthetic rectangle inside that box:

\[
x_{left} = c_x - \rho \frac{w}{2},
\qquad
x_{right} = c_x + \rho \frac{w}{2},
\]

\[
y_{top} = c_y - \rho \frac{h}{2},
\qquad
y_{bottom} = c_y + \rho \frac{h}{2},
\]

where $\rho$ is \param{light_ratio}. The four image points are:

\[
(x_{left}, y_{bottom}),\quad
(x_{left}, y_{top}),\quad
(x_{right}, y_{top}),\quad
(x_{right}, y_{bottom}).
\]

The 3D object model is a flat armor plate. The solver tries both dimensions:

\begin{itemize}
  \item small armor: $140$ mm by $125$ mm,
  \item large armor: $235$ mm by $127$ mm.
\end{itemize}

It keeps the model with lower reprojection error. Reprojection error means:

\begin{enumerate}
  \item Estimate the pose.
  \item Project the known 3D corners back into the image.
  \item Measure the average pixel distance between projected corners and input
  corners.
\end{enumerate}

The code rejects a PnP result when the mean reprojection error is too large.
The debug topic also reports a normalized error:

\[
e_{norm} = \frac{e_{pixels}}{\sqrt{w^2+h^2}}.
\]

This makes errors easier to compare across targets of different pixel size.

\section{Armor Yaw And Obliquity}

The PnP solver estimates both position and orientation. Orientation can flip
for flat rectangles because two poses can project similarly into the image.
The frame transformer protects the tracker from a bad plate normal.

Let the transformed armor center be $(x,y,z)$ in odom. The bearing from the
robot to the armor is:

\[
\theta_{bearing} = \operatorname{atan2}(y,x).
\]

If the PnP yaw points the wrong way, the code adds $\pi$. If the yaw is still
too far from the bearing by more than \param{max_oblique_deg}, the code clamps
the yaw to the bearing. This means the position can still be used even when the
plate orientation is unreliable.

\section{Tracker State}

The tracker models the enemy robot as a rotating body with armor plates on a
circle around the robot center. The EKF state vector has nine values:

\[
\vect{x} =
\begin{bmatrix}
x_c & v_{xc} & y_c & v_{yc} & z_a & v_{za} & \psi & \dot{\psi} & r
\end{bmatrix}^T.
\]

\begin{longtable}{p{0.15\textwidth}p{0.65\textwidth}p{0.12\textwidth}}
\toprule
Symbol & Meaning & Unit \\
\midrule
$x_c,y_c$ & enemy robot center position in odom & m \\
$v_{xc},v_{yc}$ & enemy robot center horizontal velocity & m/s \\
$z_a$ & armor height & m \\
$v_{za}$ & armor vertical velocity & m/s \\
$\psi$ & visible armor face yaw & rad \\
$\dot{\psi}$ & armor yaw rate & rad/s \\
$r$ & radius from enemy center to armor plate & m \\
\bottomrule
\end{longtable}

The measurement vector is:

\[
\vect{z} =
\begin{bmatrix}
x_a & y_a & z_a & \psi
\end{bmatrix}^T,
\]

where $(x_a,y_a,z_a)$ is the detected armor center in odom.

\section{Measurement Model}

The tracker assumes the armor is offset from the enemy center by radius $r$:

\[
x_a = x_c - r\cos\psi,
\qquad
y_a = y_c - r\sin\psi,
\qquad
z_a = z_a,
\qquad
\psi = \psi.
\]

In EKF notation this is:

\[
\vect{z}_{pred} = h(\vect{x}).
\]

The innovation is the difference between what was measured and what the
current state predicted:

\[
\vect{y} = \vect{z} - h(\vect{x}).
\]

A small innovation means the new detection agrees with the tracker. A large
innovation means the detection may be noisy, a different target, or a different
armor face after the enemy rotated.

\section{Why The Tracker Is An EKF}

A normal Kalman filter works when the equations are linear. This tracker has
$\sin\psi$ and $\cos\psi$, so the measurement model is nonlinear. An Extended
Kalman Filter solves this by linearizing the nonlinear equation near the
current estimate.

The linearization is the Jacobian matrix $H$:

\[
H =
\frac{\partial h}{\partial \vect{x}}.
\]

For the first two measurement rows:

\[
\frac{\partial x_a}{\partial x_c}=1,\quad
\frac{\partial x_a}{\partial \psi}=r\sin\psi,\quad
\frac{\partial x_a}{\partial r}=-\cos\psi,
\]

\[
\frac{\partial y_a}{\partial y_c}=1,\quad
\frac{\partial y_a}{\partial \psi}=-r\cos\psi,\quad
\frac{\partial y_a}{\partial r}=-\sin\psi.
\]

These are exactly the non-zero entries built in
\file{src/src/tracker.cpp} inside \code{Tracker::ekfUpdate}.

\section{Prediction Step}

Between frames, the target is predicted with a damped constant-velocity model.
For a position and velocity pair:

\[
p_{new} = p + b v \Delta t,
\qquad
v_{new} = b v.
\]

The damping coefficient $b$ is computed from \param{alpha_pos},
\param{alpha_yaw}, \param{alpha_coast}, the frame time $\Delta t$, and
\param{ref_freq}. This keeps damping approximately consistent even if the
detector rate changes.

The covariance matrix $P$ stores uncertainty. During prediction it becomes:

\[
P_{new} = FPF^T + Q.
\]

$F$ is the Jacobian of the motion model. $Q$ is process noise. Larger
\param{q_pos} or \param{q_yaw} makes the tracker respond faster but jitter
more. Smaller values make it smoother but slower.

\section{Update Step}

When a detection is associated with the current track, the Kalman update is:

\[
S = HPH^T + R,
\]

\[
K = PH^T S^{-1},
\]

\[
\vect{x}_{new} = \vect{x} + K\vect{y}.
\]

$R$ is measurement noise. In this repository, $R$ grows with target distance.
The code also increases yaw and position noise when the armor is viewed
obliquely, because a nearly edge-on plate gives weaker orientation information.

The covariance update uses Joseph form:

\[
P_{new} = (I-KH)P(I-KH)^T + KRK^T.
\]

Joseph form costs a little more computation but is more numerically stable.

\section{Mahalanobis Distance}

The tracker must decide whether a detection belongs to the current target.
It uses Mahalanobis distance:

\[
d^2 = \vect{y}^T S^{-1} \vect{y}.
\]

This is better than plain Euclidean distance because it accounts for
uncertainty. A detection far from a very uncertain prediction may still be
reasonable. A detection close to a very certain prediction can be suspicious if
the covariance says it should not move there.

The parameter \param{maha_threshold} is the gate. If the distance is above
this threshold, the detection is not associated with the current track.

\section{Tracker State Machine}

The tracker has four states:

\begin{longtable}{p{0.22\textwidth}p{0.68\textwidth}}
\toprule
State & Meaning \\
\midrule
\code{LOST} & No target is trusted. The next valid detection initializes a new
track. \\
\code{DETECTING} & A possible target has been seen but needs consecutive matches
before commands are trusted. \\
\code{TRACKING} & The target is locked. The node may publish nonzero yaw/pitch
and the fire gate may allow shooting. \\
\code{TEMP_LOST} & The tracker had a target but missed recent detections. It
coasts briefly before becoming \code{LOST}. \\
\bottomrule
\end{longtable}

The purpose of this state machine is to avoid commanding the gimbal from a
single accidental detection.

\section{Armor Face Switching}

When an enemy spins, the visible armor face changes. The tracker detects a
large yaw jump. A jump near $90^\circ$ means the visible face changed to the
alternate pair. The code swaps \code{radius_} and \code{other_radius_}, updates
the height offset \code{dz_}, and then performs a normal EKF update from the
new face.

This is why the tracker state includes the enemy center and radius instead of
only tracking the visible armor point. The center lets the code predict all
four plates.

\section{Aim Planning}

The aim planner predicts four possible future armor faces:

\[
\psi_i(t) = \psi + \dot{\psi}t + i\frac{\pi}{2},
\qquad i \in \{0,1,2,3\}.
\]

The enemy center is also predicted:

\[
x_c(t) = x_c + v_{xc}t,\qquad
y_c(t) = y_c + v_{yc}t,\qquad
z_c(t) = z_a + v_{za}t.
\]

For face $i$, the face center is:

\[
x_i = x_c(t) - r_i\cos\psi_i(t),
\qquad
y_i = y_c(t) - r_i\sin\psi_i(t),
\qquad
z_i = z_c(t) + dz_i.
\]

The planner evaluates all four faces and chooses the best one. The first pass
uses a rough prediction time:

\[
t_{pred} = t_{bias} + \frac{1.5}{v_{bullet}},
\]

where $t_{bias}$ is the \param{time_bias} parameter, or measured latency when
that feature is enabled. The second pass reuses the best face's computed
bullet flight time.

\section{Barrel Origin And Parallax}

The aim vector starts at the barrel, not at the camera. This matters at close
range. A $10$ cm camera-to-barrel offset at $1$ m is a large angular error.

The barrel offset is measured in the gimbal body frame:

\[
\vect{o}_{barrel} =
\begin{bmatrix}
o_x\\o_y\\o_z
\end{bmatrix}.
\]

The code rotates horizontal offset by current yaw:

\[
b_x = o_x\cos\theta - o_y\sin\theta,
\qquad
b_y = o_x\sin\theta + o_y\cos\theta,
\qquad
b_z = h_{gimbal} + o_z.
\]

Then the aim vector is face position minus barrel position.

\section{Ballistics}

Without gravity, pitch would be:

\[
\phi = \operatorname{atan2}(\Delta z, d),
\]

where $d$ is horizontal distance and $\Delta z$ is vertical target difference.
Gravity makes the bullet drop during flight, so the barrel must aim higher.

The solver iterates:

\[
v_x = v\cos\phi,
\qquad
t = \frac{d}{v_x},
\]

\[
\phi = \operatorname{atan2}\left(\Delta z + \frac{1}{2}gt^2, d\right).
\]

After five iterations, it returns pitch and flight time. It rejects impossible
geometry such as nearly vertical shots or a pitch magnitude above about
$1.2$ radians.

\section{Fire Margin}

For each predicted face, the planner computes an angular window:

\[
w = w_0 \min\left(\frac{d_{ref}}{\max(R,0.5)},2.0\right).
\]

Here $w_0$ is \param{angular_window}, $d_{ref}$ is
\param{window_ref_dist}, and $R$ is target range.

The error is:

\[
e = |\operatorname{shortest\_angular\_distance}(bearing, face\_yaw)|.
\]

The margin is:

\[
margin = w - e.
\]

If the margin is positive, the selected face is aligned enough according to
the planner. The fire gate still performs more checks before the actual fire
flag is published.

\section{Anti-Gyro Timing}

When anti-gyro mode is enabled and the enemy yaw rate is high, face alignment
timing matters. The residual is:

\[
residual = \frac{
\operatorname{shortest\_angular\_distance}(face\_yaw,bearing)
}{\dot{\psi}}.
\]

Its unit is seconds. A residual near zero means the bullet impact time and the
face alignment time match. The fire gate can block shots when
$|residual|$ is larger than \param{fire.anti_gyro_max_residual}.

\section{Command Smoothing}

The node smooths yaw and pitch commands using an exponential moving average.
Pitch uses:

\[
p_{smooth} = \alpha p_{target} + (1-\alpha)p_{old}.
\]

Yaw uses the shortest angular distance so that a target crossing from
$+\pi$ to $-\pi$ does not command a full revolution:

\[
y_{smooth} =
normalize(y_{old} + \alpha \cdot shortest(y_{old}, y_{target})).
\]

\param{cmd_smooth_alpha}=1 means no smoothing. Smaller values smooth more but
increase lag. The fire gate blocks shooting when smoothing lag is too large.

\section{Fire Gate Logic}

The fire gate is intentionally conservative. It allows fire only when all
required checks pass:

\begin{enumerate}
  \item Tracker is in \code{TRACKING}.
  \item Aim target is valid.
  \item Ballistic solution is valid.
  \item Distance is within \param{min_fire_dist} and \param{max_fire_dist}.
  \item Optional stale-measurement and anti-gyro checks pass.
  \item Planner margin is non-negative.
  \item Camera-frame aim angle is inside \param{angular_window}.
  \item Smoothed command is close enough to the unsmoothed target.
\end{enumerate}

If a margin-like check fails briefly after a shot was allowed, hysteresis can
keep fire enabled across small angular drift. Hysteresis is dropped as soon as
a hard check fails.

\section{Debugging With The Math}

When the robot misses or jitters, debug in the same order as the math pipeline:

\begin{enumerate}
  \item Detection: Are bounding boxes stable in pixels?
  \item PnP: Is reprojection error low and depth stable?
  \item Transform: Is odom position stable when the robot is still?
  \item EKF: Are innovation and Mahalanobis distance reasonable?
  \item Aim: Is the selected face stable and the predicted target plausible?
  \item Command: Is smoothing lag large?
  \item Fire gate: Which blocker reason dominates?
\end{enumerate}

Do not start by changing EKF noise if PnP depth is already noisy. Do not tune
ballistics if the measured bullet speed is unknown. Fix the earliest wrong
stage.

\section{Practice Problems}

\begin{enumerate}
  \item A detector reports a box centered at $(500,300)$ with width $100$ px,
  height $80$ px, and \param{light_ratio}=0.85. Compute the four synthetic PnP
  image points.
  \item A point has camera coordinates $(0.2, -0.1, 3.0)$ m. With
  $f_x=f_y=700$, $c_x=480$, $c_y=300$, compute the projected pixel.
  \item If \param{bullet_speed} is accidentally set too high, will long-range
  shots tend to hit high or low? Explain using the gravity equation.
  \item A target has $\operatorname{shortest}(bearing, face\_yaw)=0.06$ rad and
  the scaled window is $0.05$ rad. What is the margin? Can the planner allow
  fire?
  \item Why is Mahalanobis distance better than plain distance for target
  association?
\end{enumerate}
